% !PRINTED 139
\weberpage{139}

These functions satisfy, by (9), the differential equations
\weq{11}{\begin{aligned}
\frac{dx}{dv} &= yz,\\
\frac{dy}{dv} &= -zx,\\
\frac{dz}{dv} &= -\varkappa^2xy,
\end{aligned}}
and the subsidiary conditions:
\weq{12}{\operatorname{sn}0 = 0, \qquad \operatorname{cn}0 = 1, \qquad \operatorname{dn}0 = 1.}

Between them there subsist the relations
\weq{13}{y^2 = 1 - x^2, \qquad z^2 = 1 - \varkappa^2x^2.}

From the fundamental properties of the $\vartheta$-functions there follow
the first properties of the elliptic functions:
\[
\operatorname{sn}v = -\operatorname{sn}(-v), \qquad \operatorname{cn}v = \operatorname{cn}(-v), \qquad \operatorname{dn}v = \operatorname{dn}(-v),
\]
that is, $\operatorname{sn}v$ is an odd function, while $\operatorname{cn}v$, $\operatorname{dn}v$ are even functions.

If we put further
\weq{14}{\pi\vartheta_{00}^2 = 2K, \qquad \pi\vartheta_{00}^2\omega = 2iK',}
and consequently
\weq{15}{\pi\vartheta_{01}^2 = 2\varkappa'K, \qquad \pi\vartheta_{10}^2 = 2\varkappa K, \qquad \omega = \frac{iK'}{K},}
then $v$ takes the values $K$, $iK'$, $K+iK'$ when $u = \frac{1}{2}, \frac{\omega}{2}, \frac{1+\omega}{2}$, and it follows from (4), having regard to (3) and \S\,21, (10):
\weq{16}{\begin{aligned}
&\operatorname{sn}K = 1, \quad \operatorname{sn}(K+iK') = \frac{1}{\varkappa},\\
&\operatorname{cn}K = 0, \quad \operatorname{cn}(K+iK') = -\frac{i\varkappa'}{\varkappa},\\
&\operatorname{dn}K = \varkappa', \quad \operatorname{dn}(K+iK') = 0,\\
&\operatorname{sn}iK', \quad \operatorname{cn}iK', \quad \operatorname{dn}iK' = \infty.
\end{aligned}}

Now $K$, $K'$ can be expressed by means of (8) in terms of definite
integrals, and one obtains, if one displaces $v$ from one vertex of the
parallelogram $0, K, K+iK', iK'$ to the next along the periphery, and
hence $u$ along
\[
0, \qquad \frac{1}{2}, \qquad \frac{1+\omega}{2}, \qquad \frac{\omega}{2}.
\]
